% Options for packages loaded elsewhere
\PassOptionsToPackage{unicode}{hyperref}
\PassOptionsToPackage{hyphens}{url}
%
\documentclass[
]{article}
\usepackage{amsmath,amssymb}
\usepackage{iftex}
\ifPDFTeX
  \usepackage[T1]{fontenc}
  \usepackage[utf8]{inputenc}
  \usepackage{textcomp} % provide euro and other symbols
\else % if luatex or xetex
  \usepackage{unicode-math} % this also loads fontspec
  \defaultfontfeatures{Scale=MatchLowercase}
  \defaultfontfeatures[\rmfamily]{Ligatures=TeX,Scale=1}
\fi
\usepackage{lmodern}
\ifPDFTeX\else
  % xetex/luatex font selection
\fi
% Use upquote if available, for straight quotes in verbatim environments
\IfFileExists{upquote.sty}{\usepackage{upquote}}{}
\IfFileExists{microtype.sty}{% use microtype if available
  \usepackage[]{microtype}
  \UseMicrotypeSet[protrusion]{basicmath} % disable protrusion for tt fonts
}{}
\makeatletter
\@ifundefined{KOMAClassName}{% if non-KOMA class
  \IfFileExists{parskip.sty}{%
    \usepackage{parskip}
  }{% else
    \setlength{\parindent}{0pt}
    \setlength{\parskip}{6pt plus 2pt minus 1pt}}
}{% if KOMA class
  \KOMAoptions{parskip=half}}
\makeatother
\usepackage{xcolor}
\usepackage[margin=1in]{geometry}
\usepackage{color}
\usepackage{fancyvrb}
\newcommand{\VerbBar}{|}
\newcommand{\VERB}{\Verb[commandchars=\\\{\}]}
\DefineVerbatimEnvironment{Highlighting}{Verbatim}{commandchars=\\\{\}}
% Add ',fontsize=\small' for more characters per line
\usepackage{framed}
\definecolor{shadecolor}{RGB}{248,248,248}
\newenvironment{Shaded}{\begin{snugshade}}{\end{snugshade}}
\newcommand{\AlertTok}[1]{\textcolor[rgb]{0.94,0.16,0.16}{#1}}
\newcommand{\AnnotationTok}[1]{\textcolor[rgb]{0.56,0.35,0.01}{\textbf{\textit{#1}}}}
\newcommand{\AttributeTok}[1]{\textcolor[rgb]{0.13,0.29,0.53}{#1}}
\newcommand{\BaseNTok}[1]{\textcolor[rgb]{0.00,0.00,0.81}{#1}}
\newcommand{\BuiltInTok}[1]{#1}
\newcommand{\CharTok}[1]{\textcolor[rgb]{0.31,0.60,0.02}{#1}}
\newcommand{\CommentTok}[1]{\textcolor[rgb]{0.56,0.35,0.01}{\textit{#1}}}
\newcommand{\CommentVarTok}[1]{\textcolor[rgb]{0.56,0.35,0.01}{\textbf{\textit{#1}}}}
\newcommand{\ConstantTok}[1]{\textcolor[rgb]{0.56,0.35,0.01}{#1}}
\newcommand{\ControlFlowTok}[1]{\textcolor[rgb]{0.13,0.29,0.53}{\textbf{#1}}}
\newcommand{\DataTypeTok}[1]{\textcolor[rgb]{0.13,0.29,0.53}{#1}}
\newcommand{\DecValTok}[1]{\textcolor[rgb]{0.00,0.00,0.81}{#1}}
\newcommand{\DocumentationTok}[1]{\textcolor[rgb]{0.56,0.35,0.01}{\textbf{\textit{#1}}}}
\newcommand{\ErrorTok}[1]{\textcolor[rgb]{0.64,0.00,0.00}{\textbf{#1}}}
\newcommand{\ExtensionTok}[1]{#1}
\newcommand{\FloatTok}[1]{\textcolor[rgb]{0.00,0.00,0.81}{#1}}
\newcommand{\FunctionTok}[1]{\textcolor[rgb]{0.13,0.29,0.53}{\textbf{#1}}}
\newcommand{\ImportTok}[1]{#1}
\newcommand{\InformationTok}[1]{\textcolor[rgb]{0.56,0.35,0.01}{\textbf{\textit{#1}}}}
\newcommand{\KeywordTok}[1]{\textcolor[rgb]{0.13,0.29,0.53}{\textbf{#1}}}
\newcommand{\NormalTok}[1]{#1}
\newcommand{\OperatorTok}[1]{\textcolor[rgb]{0.81,0.36,0.00}{\textbf{#1}}}
\newcommand{\OtherTok}[1]{\textcolor[rgb]{0.56,0.35,0.01}{#1}}
\newcommand{\PreprocessorTok}[1]{\textcolor[rgb]{0.56,0.35,0.01}{\textit{#1}}}
\newcommand{\RegionMarkerTok}[1]{#1}
\newcommand{\SpecialCharTok}[1]{\textcolor[rgb]{0.81,0.36,0.00}{\textbf{#1}}}
\newcommand{\SpecialStringTok}[1]{\textcolor[rgb]{0.31,0.60,0.02}{#1}}
\newcommand{\StringTok}[1]{\textcolor[rgb]{0.31,0.60,0.02}{#1}}
\newcommand{\VariableTok}[1]{\textcolor[rgb]{0.00,0.00,0.00}{#1}}
\newcommand{\VerbatimStringTok}[1]{\textcolor[rgb]{0.31,0.60,0.02}{#1}}
\newcommand{\WarningTok}[1]{\textcolor[rgb]{0.56,0.35,0.01}{\textbf{\textit{#1}}}}
\usepackage{graphicx}
\makeatletter
\newsavebox\pandoc@box
\newcommand*\pandocbounded[1]{% scales image to fit in text height/width
  \sbox\pandoc@box{#1}%
  \Gscale@div\@tempa{\textheight}{\dimexpr\ht\pandoc@box+\dp\pandoc@box\relax}%
  \Gscale@div\@tempb{\linewidth}{\wd\pandoc@box}%
  \ifdim\@tempb\p@<\@tempa\p@\let\@tempa\@tempb\fi% select the smaller of both
  \ifdim\@tempa\p@<\p@\scalebox{\@tempa}{\usebox\pandoc@box}%
  \else\usebox{\pandoc@box}%
  \fi%
}
% Set default figure placement to htbp
\def\fps@figure{htbp}
\makeatother
\setlength{\emergencystretch}{3em} % prevent overfull lines
\providecommand{\tightlist}{%
  \setlength{\itemsep}{0pt}\setlength{\parskip}{0pt}}
\setcounter{secnumdepth}{-\maxdimen} % remove section numbering
\usepackage{bookmark}
\IfFileExists{xurl.sty}{\usepackage{xurl}}{} % add URL line breaks if available
\urlstyle{same}
\hypersetup{
  pdftitle={1-Solutions-GenDistance},
  hidelinks,
  pdfcreator={LaTeX via pandoc}}

\title{1-Solutions-GenDistance}
\author{}
\date{\vspace{-2.5em}2026-08-31}

\begin{document}
\maketitle

\section{Part I: Read in and Align the Sequences with
Clustal-Omega}\label{part-i-read-in-and-align-the-sequences-with-clustal-omega}

\begin{Shaded}
\begin{Highlighting}[]
\CommentTok{\#Align sequences}
\NormalTok{dnaSeqs }\OtherTok{=} \FunctionTok{readDNAStringSet}\NormalTok{(}\StringTok{"flu\_seqs.fasta"}\NormalTok{)}
\NormalTok{aligns }\OtherTok{=} \FunctionTok{msa}\NormalTok{(dnaSeqs, }\AttributeTok{method=}\StringTok{"ClustalOmega"}\NormalTok{, }\AttributeTok{order=}\StringTok{"input"}\NormalTok{)}
\end{Highlighting}
\end{Shaded}

\begin{verbatim}
## using Gonnet
\end{verbatim}

\begin{Shaded}
\begin{Highlighting}[]
\CommentTok{\# ── Load \& normalize ──────────────────────────────────────────────────────────}
\end{Highlighting}
\end{Shaded}

\section{Part II: Use the Phangorn package to compare substitution
models}\label{part-ii-use-the-phangorn-package-to-compare-substitution-models}

forPhang = msaConvert(aligns, type=``phangorn::phyDat'')

mt = modelTest(forPhang,
model=c(``JC'',``F81'',``K80'',``HKY'',``TrN'')) write.table(mt,
``modelTest\_results.txt'', quote=FALSE, sep=``\t",row.names=FALSE,
col.names=TRUE)

\begin{verbatim}

## Question 1: The Bayesian Information Criterion (BIC) is a method for comparing models that includes penalties for more complex models. With this method, the best one is the model with the LOWEST BIC score. In your test, which model is this? Describe this model in terms of its assumptions. (5 pts)

YOUR ANSWER TO QUESTION ONE HERE.

# Part III: Calculate Genetic Distance Under the Best Model

## Question 2: Show your command for calculating distance, including which model you used in your command, and which option you used for gamma. (1 pt)


``` r
#->YOUR CODE HERE
\end{verbatim}

\section{Part IV: Create both UPGMA and Neighbor-Joining
Trees}\label{part-iv-create-both-upgma-and-neighbor-joining-trees}

\begin{Shaded}
\begin{Highlighting}[]
\CommentTok{\#{-}\textgreater{}YOUR CODE HERE}
\end{Highlighting}
\end{Shaded}

\subsection{Question 3: Provide your plots for the rooted UPGMA and
rooted Neighbor-Joining trees (be sure they are labeled so I can tell
which is which!). Describe any differences you see in the overall
appearance of the two trees. Do the same clades appear in each tree? Do
viruses collected in later years branch off from the viruses of previous
years (like we would expect)? (5
pts)}\label{question-3-provide-your-plots-for-the-rooted-upgma-and-rooted-neighbor-joining-trees-be-sure-they-are-labeled-so-i-can-tell-which-is-which.-describe-any-differences-you-see-in-the-overall-appearance-of-the-two-trees.-do-the-same-clades-appear-in-each-tree-do-viruses-collected-in-later-years-branch-off-from-the-viruses-of-previous-years-like-we-would-expect-5-pts}

YOUR ANSWER TO QUESTION 3 HERE.

\section{Part V: Calculate Distance Metrics Between the 2
Trees}\label{part-v-calculate-distance-metrics-between-the-2-trees}

\subsection{Question 4: What is the symmetric difference, branch score
difference, and path difference between your two trees? Which number
implies the biggest difference in the trees? Based on what you know
about the definition of these metrics, why would one score be much
higher than the others? Does your highest score make sense given what
you see when you plot the two trees? (3
pts)}\label{question-4-what-is-the-symmetric-difference-branch-score-difference-and-path-difference-between-your-two-trees-which-number-implies-the-biggest-difference-in-the-trees-based-on-what-you-know-about-the-definition-of-these-metrics-why-would-one-score-be-much-higher-than-the-others-does-your-highest-score-make-sense-given-what-you-see-when-you-plot-the-two-trees-3-pts}

YOUR ANSWER TO QUESTION 4.

\section{Part VI: Perform Bootstrapping on the Neighbor-Joining
Tree}\label{part-vi-perform-bootstrapping-on-the-neighbor-joining-tree}

\subsection{Question 5: Provide a plot of your bootstrapped tree. Are
there any nodes in your tree that seem to have very weak bootstrap
support (i.e.~less than 50\%)? How many of these weakly supported nodes
do you see? Which node has the weakest support, and what clade is this
node at the base of? If it is hard to see exactly where the nodes are
with the bootstrap labels on them, you might also want to refer back to
your original plot without the labels to help you. (3
pts)}\label{question-5-provide-a-plot-of-your-bootstrapped-tree.-are-there-any-nodes-in-your-tree-that-seem-to-have-very-weak-bootstrap-support-i.e.-less-than-50-how-many-of-these-weakly-supported-nodes-do-you-see-which-node-has-the-weakest-support-and-what-clade-is-this-node-at-the-base-of-if-it-is-hard-to-see-exactly-where-the-nodes-are-with-the-bootstrap-labels-on-them-you-might-also-want-to-refer-back-to-your-original-plot-without-the-labels-to-help-you.-3-pts}

\begin{Shaded}
\begin{Highlighting}[]
\CommentTok{\#{-}\textgreater{}YOUR CODE HERE}
\end{Highlighting}
\end{Shaded}

YOUR ANSWER TO QUESTION 5.

\subsection{Question 6: Provide a plot of your majority rule consensus
tree. (1
pt)}\label{question-6-provide-a-plot-of-your-majority-rule-consensus-tree.-1-pt}

\begin{Shaded}
\begin{Highlighting}[]
\CommentTok{\#{-}\textgreater{}YOUR CODE HERE}
\end{Highlighting}
\end{Shaded}

\section{Part VII: Find the Most Parsimonious
Tree}\label{part-vii-find-the-most-parsimonious-tree}

\subsection{Question 7: What are your parsimony scores for your UPGMA
and Neighbor-Joining trees? Given that a lower score is better, which of
your trees is the most parsimonious? Based on what you know about the
definition of parsimony, what does it mean when we say that one tree is
more parsimonious than another? i.e.~How should we interpret this
result? (3
pts)}\label{question-7-what-are-your-parsimony-scores-for-your-upgma-and-neighbor-joining-trees-given-that-a-lower-score-is-better-which-of-your-trees-is-the-most-parsimonious-based-on-what-you-know-about-the-definition-of-parsimony-what-does-it-mean-when-we-say-that-one-tree-is-more-parsimonious-than-another-i.e.-how-should-we-interpret-this-result-3-pts}

\begin{Shaded}
\begin{Highlighting}[]
\CommentTok{\#{-}\textgreater{}YOUR CODE HERE}
\end{Highlighting}
\end{Shaded}

YOUR ANSWER TO QUESTION 7.

\subsection{Question 8: Were you able to improve on the parsimony score
you got with the neighbor-joining tree? (1
pt)}\label{question-8-were-you-able-to-improve-on-the-parsimony-score-you-got-with-the-neighbor-joining-tree-1-pt}

YOUR ANSWER TO QUESTION 8.

\subsection{Question 9: Provide the plot for your final tree. Even
though there are still some uncertain relationships, can you now see any
kind of trend with respect to the years the viruses were sampled and
where they appear on the tree? Describe your new tree and the trends
that you see. (3
pts)}\label{question-9-provide-the-plot-for-your-final-tree.-even-though-there-are-still-some-uncertain-relationships-can-you-now-see-any-kind-of-trend-with-respect-to-the-years-the-viruses-were-sampled-and-where-they-appear-on-the-tree-describe-your-new-tree-and-the-trends-that-you-see.-3-pts}

\begin{Shaded}
\begin{Highlighting}[]
\CommentTok{\#{-}\textgreater{}YOUR CODE HERE}
\end{Highlighting}
\end{Shaded}

YOUR ANSWER TO QUESTION 9.

\end{document}
